\section{Old-label certification through one repair}\label{sec:certification}
\subsection{The ordered experiment}
We now impose a narrower experiment than the control model above. Let
$\theta_i\in\{0,1\}$ denote adequacy of the \emph{old} representation.
It remains the terminal target after repair.

\begin{assumption}[Finite ordered repair experiment]\label{ass:repair}
The following parameters are fixed as $\delta\downarrow0$.
\begin{enumerate}
\item There are finitely many environments. Pre-repair observations are iid
with pairwise distinct laws $P_i$ on a common finite alphabet, all with full
support. The binary labels include both values.
\item A pre-repair task slot costs $g_i>0$ and can supply one observation from
$P_i$. There is no free informative pre-repair action. Pausing or completing
old-state tasks does not remove their cost.
\item After entering the repair phase, no more pre-repair observations are
available. Each attempt takes one slot and costs $\kappa_i\ge0$, including its
opportunity cost. It succeeds with known probability $q>0$, independent of
the environment, previous observations, and other attempts. Success is
observable and irreversible. A failed attempt can be retried.
\item After success, each slot supplies an iid observation with law $Q_i$
on a finite alphabet at zero cost. These observations are independent of
the pre-repair stream and repair trials conditional on $i$. Distinct $Q_i$
need not exist for every pair. For simplicity all $Q_i$ have full support.
\item All action choices use observed data and independent policy randomization.
Accounting costs are not additional signals. All slots up to the prescribed
deadline are accounted for, and a declaration is produced by that deadline.
\end{enumerate}
\end{assumption}
The cost $\kappa_i$ is the total cost of a repair slot, not an extra fee to
which an omitted old-state loss must later be added. Post-repair observations
are free in cost, not in time. Policies may stop trying repair, in which case
remaining old-state task slots incur $g_i$. A policy that never repairs has
only a direct pre-repair declaration. These conventions exclude uncharged
tails while retaining the natural direct-certification alternative.

Let $L_\delta=\log(1/\delta)$ and
\begin{equation}\label{eq:horizon}
 H_\delta=\lceil L_\delta^2\rceil.
\end{equation}
For sufficiently small $\delta$, let $\Pi_\delta$ be the ordered policies
with $\PP_i^\pi(\widehat\theta\ne\theta_i)\le\delta$ for every $i$, and define
\[
 J_i^*(\delta)=\inf_{\pi\in\Pi_\delta}\EE_i^\pi C.
\]
The separated class $\Pi_\delta^{\rm sep}$ additionally commits to its final
old-label answer \emph{before its first repair attempt}, using only pre-repair
observations and policy coins; it must meet the same uniform error guarantee.
A no-repair policy commits by the deadline. Define $J_i^{\rm sep}(\delta)$
by the corresponding infimum. This is an operational certification-before-
repair restriction, not a Bayes posterior threshold.

\subsection{Post-repair classes and admissible answers}
Let $\calC$ be the partition induced by exact equality of the post-repair
laws, and write $C_i$ for the class containing $i$. A decoder is a map
$a:\calC\to\{0,1\}$. Define
\begin{equation}\label{eq:answers}
 \cD_i=\{a:a(C_i)=\theta_i\},\qquad
 B(a)=\{j:a(C_j)\ne\theta_j\}.
\end{equation}
The decoder is valid in environment $i$ exactly when $a\in\cD_i$.
Its invalidating environments may have either old label.
For a valid decoder, set
\begin{align}
 d_i(a)&=\min_{j\in B(a)}\KL(P_i\Vert P_j),\\
 D_i&=\max_{a\in\cD_i}d_i(a),&
 D_i^{\rm sep}&=\min_{j:\theta_j\ne\theta_i}\KL(P_i\Vert P_j).
 \label{eq:rates}
\end{align}
There are only $2^{|\calC|}$ decoders. No optimization over observation
allocation appears because there is exactly one pre-repair channel.

\begin{theorem}[Decoder cost coefficient]\label{thm:main}
Under Assumption~\ref{ass:repair}, as $\delta\downarrow0$,
\begin{equation}\label{eq:main}
 \frac{J_i^*(\delta)}{L_\delta}\longrightarrow\frac{g_i}{D_i},
 \qquad
 \frac{J_i^{\rm sep}(\delta)}{L_\delta}\longrightarrow
 \frac{g_i}{D_i^{\rm sep}}\quad\text{for every }i.
\end{equation}
A single sequence of implementable joint policies attains the joint
coefficients simultaneously in all environments. If a decoder is valid in
every environment, the joint coefficient is zero and a policy of uniformly
bounded expected cost exists for all sufficiently small $\delta$.
\end{theorem}
Because the constant decoder $a(C)=\theta_i$ is valid in $i$,
$D_i\ge D_i^{\rm sep}$ and joint repair cannot have a larger leading
coefficient. This comparison does not say that a finite-$\delta$ cost ratio
equals the asymptotic coefficient ratio.

\subsection{Three environments and same-label exclusion}\label{sec:example}
Take Bernoulli laws with
\begin{equation}\label{eq:example}
 (p_1,p_2,p_3)=(0.5,0.9,0.6),\quad
 (q_1,q_2,q_3)=(0.8,0.2,0.2),\quad
 (\theta_1,\theta_2,\theta_3)=(0,0,1),\quad g_i=1.
\end{equation}
Here $q_i$ denotes a post-repair mean, whereas $q$ without a subscript denotes
the independent repair success probability.
The classes are $C_A=\{1\}$ and $C_B=\{2,3\}$.
Listing decoders by $(a(C_A),a(C_B))$ gives
\begin{center}
\begin{tabular}{ccl}
\toprule
Decoder & $B(a)$ & Valid environments\\
\midrule
$(0,0)$ & $\{3\}$ & $1,2$\\
$(0,1)$ & $\{2\}$ & $1,3$\\
$(1,0)$ & $\{1,3\}$ & $2$\\
$(1,1)$ & $\{1,2\}$ & $3$\\
\bottomrule
\end{tabular}
\end{center}
Under $E_1$, the best decoder is $(0,1)$. It excludes $E_2$, which shares
$E_1$'s old label, while retaining the opposite-label $E_3$ for post-repair
discrimination. Direct old-label certification instead must exclude $E_3$.
For Bernoulli means,
\[
 \KL(\Bern(p)\Vert\Bern(r))
 =p\log(p/r)+(1-p)\log((1-p)/(1-r)).
\]
Thus
\begin{align*}
D_1&=\KL(\Bern(0.5)\Vert\Bern(0.9))=0.510825623766\ldots,\\
D_1^{\rm sep}&=\KL(\Bern(0.5)\Vert\Bern(0.6))=0.020410997260\ldots,\\
\frac{g_1}{D_1}&=1.957615188971\ldots,\qquad
\frac{g_1}{D_1^{\rm sep}}=48.993196523204\ldots,\\
\frac{J_1^{\rm sep}(\delta)}{J_1^*(\delta)}
 &\longrightarrow 25.026980174255\ldots.
\end{align*}
The other environments and their best decoders are calculated in
Table~\ref{tab:analytic}. These values are generated from the model file.

\begin{table}[ht]
\centering\small\begin{tabular}{lrrrr}
\toprule
Environment & $D_i$ & $g_i/D_i$ & $g_i/D_i^{\rm sep}$ & Decoder\\
\midrule
$E_1$ & 0.510826 & 1.957615 & 48.993197 & $(0,1)$\\
$E_2$ & 0.226289 & 4.419125 & 4.419125 & $(0,0)$\\
$E_3$ & 0.311239 & 3.212968 & 49.663496 & $(0,1)$\\
\bottomrule
\end{tabular}

\caption{Analytic decoder coefficients for \eqref{eq:example}. Costs have
$g_i=1$. Decoder coordinates follow $(C_A,C_B)$.}\label{tab:analytic}
\end{table}

\subsection{Partition boundaries}
\begin{corollary}[Preservation, erasure, and refinement]\label{cor:partition}
Under the theorem's assumptions:
\begin{enumerate}
\item If every post-repair class has a constant old label, a universally valid
decoder exists, $D_i=+\infty$, and bounded expected joint cost is achievable.
\item If all post-repair laws are equal, $D_i=D_i^{\rm sep}$ for every $i$.
\item Refining the post-repair partition while keeping the pre-repair experiment
and labels fixed cannot decrease $D_i$.
\end{enumerate}
\end{corollary}
\begin{proof}
For the first statement, assign each class its common label. No environment
invalidates that decoder. For the second, a decoder on one class is a constant
label, so its invalidating set is exactly the opposite-label set.
For the third, lift any coarse decoder to a finer partition by assigning the
parent class's value to each child. Its invalidating set stays unchanged, while
the finer partition may allow additional decoders. Maximizing over that larger
set cannot reduce the rate.
\end{proof}
Partition equality here is exact. Two nearly equal but distinct post laws
remain separate classes in the asymptotic theorem, although distinguishing
them at practical deadlines can be expensive in time.
